\documentclass[12pt]{extarticle}
\newcommand{\level}{Grades 4--5}
\newcommand{\packetid}{F07-U-v4}
% Shared preamble for the Week 7 student pages (take-away games).
% Each packet defines \level and \packetid before % Shared preamble for the Week 7 student pages (take-away games).
% Each packet defines \level and \packetid before % Shared preamble for the Week 7 student pages (take-away games).
% Each packet defines \level and \packetid before \input{common}.
\usepackage[T1]{fontenc}
\usepackage{newcent}
\usepackage[scaled=0.92]{helvet}
\usepackage[letterpaper,top=0.95in,bottom=0.9in,left=0.5in,right=0.5in,
  headheight=18pt,headsep=0.3in,footskip=0.45in]{geometry}
\usepackage{fancyhdr}
\usepackage{tikz}

\setlength{\parindent}{0pt}
\setlength{\parskip}{0pt}
\raggedright

\pagestyle{fancy}
\fancyhf{}
\lhead{Week 7 / Take-away games / \level}
\lfoot{Bellingham Math Circle / Week 7 / \packetid}
\rfoot{\thepage}
\renewcommand{\headrulewidth}{0pt}
\renewcommand{\footrulewidth}{0pt}

% ---------- problems ----------
\newcounter{prob}
\newlength{\probgap}\setlength{\probgap}{0.35in}
% A problem never splits across pages.
\newenvironment{problem}%
  {\par\noindent\begin{minipage}[t]{\linewidth}\raggedright
   \stepcounter{prob}\textbf{Problem \theprob:}\ }%
  {\end{minipage}\par\vspace{\probgap}}

% ---------- number tracks ----------
% A snake: 0 to 10 run left to right in the first row, 11 sits under 10
% and 11-20 run right to left, 21 sits under 20 and 21-30 run left to
% right. Consecutive squares always share an edge. Thick walls separate
% rows except where the track turns.
\newlength{\sq}\setlength{\sq}{0.676in}
\newcommand{\tracknum}{\sffamily\large}
\newcommand{\track}[1]{%
  \begin{tikzpicture}[x=\sq,y=\sq,line join=miter]
    \foreach \n in {0,...,#1}{%
      \pgfmathtruncatemacro{\r}{\n==0 ? 0 : div(\n-1,10)}%
      \pgfmathtruncatemacro{\k}{\n==0 ? -1 : mod(\n-1,10)}%
      \pgfmathtruncatemacro{\odd}{mod(\r,2)}%
      \pgfmathtruncatemacro{\c}{\odd==1 ? 10-\k : \k+1}%
      \draw[line width=0.8pt] (\c,-\r) rectangle ++(1,1);
      \node[font=\tracknum,inner sep=0pt,text height=1.6ex,text depth=0pt] at (\c+0.5,-\r+0.5) {\n};
    }
    \ifnum#1>20
      \draw[line width=2.2pt] (1,0) -- (10,0);
      \draw[line width=2.2pt] (2,-1) -- (11,-1);
      \draw[line width=2.2pt] (0,1) -- (11,1) -- (11,-2) -- (1,-2) -- (1,0) -- (0,0) -- cycle;
    \else\ifnum#1>10
      \draw[line width=2.2pt] (1,0) -- (10,0);
      \draw[line width=2.2pt] (0,1) -- (11,1) -- (11,-1) -- (1,-1) -- (1,0) -- (0,0) -- cycle;
    \else
      \draw[line width=2.2pt] (0,0) rectangle (11,1);
    \fi\fi
  \end{tikzpicture}}

% ---------- piles of counters ----------
% Use inside a tikzpicture with x=1in,y=1in; arguments are the centre.
\newcommand{\counter}[2]{%
  \draw[line width=1pt,fill=black!22] ({#1},{#2}) circle (0.2in);}
.
\usepackage[T1]{fontenc}
\usepackage{newcent}
\usepackage[scaled=0.92]{helvet}
\usepackage[letterpaper,top=0.95in,bottom=0.9in,left=0.5in,right=0.5in,
  headheight=18pt,headsep=0.3in,footskip=0.45in]{geometry}
\usepackage{fancyhdr}
\usepackage{tikz}

\setlength{\parindent}{0pt}
\setlength{\parskip}{0pt}
\raggedright

\pagestyle{fancy}
\fancyhf{}
\lhead{Week 7 / Take-away games / \level}
\lfoot{Bellingham Math Circle / Week 7 / \packetid}
\rfoot{\thepage}
\renewcommand{\headrulewidth}{0pt}
\renewcommand{\footrulewidth}{0pt}

% ---------- problems ----------
\newcounter{prob}
\newlength{\probgap}\setlength{\probgap}{0.35in}
% A problem never splits across pages.
\newenvironment{problem}%
  {\par\noindent\begin{minipage}[t]{\linewidth}\raggedright
   \stepcounter{prob}\textbf{Problem \theprob:}\ }%
  {\end{minipage}\par\vspace{\probgap}}

% ---------- number tracks ----------
% A snake: 0 to 10 run left to right in the first row, 11 sits under 10
% and 11-20 run right to left, 21 sits under 20 and 21-30 run left to
% right. Consecutive squares always share an edge. Thick walls separate
% rows except where the track turns.
\newlength{\sq}\setlength{\sq}{0.676in}
\newcommand{\tracknum}{\sffamily\large}
\newcommand{\track}[1]{%
  \begin{tikzpicture}[x=\sq,y=\sq,line join=miter]
    \foreach \n in {0,...,#1}{%
      \pgfmathtruncatemacro{\r}{\n==0 ? 0 : div(\n-1,10)}%
      \pgfmathtruncatemacro{\k}{\n==0 ? -1 : mod(\n-1,10)}%
      \pgfmathtruncatemacro{\odd}{mod(\r,2)}%
      \pgfmathtruncatemacro{\c}{\odd==1 ? 10-\k : \k+1}%
      \draw[line width=0.8pt] (\c,-\r) rectangle ++(1,1);
      \node[font=\tracknum,inner sep=0pt,text height=1.6ex,text depth=0pt] at (\c+0.5,-\r+0.5) {\n};
    }
    \ifnum#1>20
      \draw[line width=2.2pt] (1,0) -- (10,0);
      \draw[line width=2.2pt] (2,-1) -- (11,-1);
      \draw[line width=2.2pt] (0,1) -- (11,1) -- (11,-2) -- (1,-2) -- (1,0) -- (0,0) -- cycle;
    \else\ifnum#1>10
      \draw[line width=2.2pt] (1,0) -- (10,0);
      \draw[line width=2.2pt] (0,1) -- (11,1) -- (11,-1) -- (1,-1) -- (1,0) -- (0,0) -- cycle;
    \else
      \draw[line width=2.2pt] (0,0) rectangle (11,1);
    \fi\fi
  \end{tikzpicture}}

% ---------- piles of counters ----------
% Use inside a tikzpicture with x=1in,y=1in; arguments are the centre.
\newcommand{\counter}[2]{%
  \draw[line width=1pt,fill=black!22] ({#1},{#2}) circle (0.2in);}
.
\usepackage[T1]{fontenc}
\usepackage{newcent}
\usepackage[scaled=0.92]{helvet}
\usepackage[letterpaper,top=0.95in,bottom=0.9in,left=0.5in,right=0.5in,
  headheight=18pt,headsep=0.3in,footskip=0.45in]{geometry}
\usepackage{fancyhdr}
\usepackage{tikz}

\setlength{\parindent}{0pt}
\setlength{\parskip}{0pt}
\raggedright

\pagestyle{fancy}
\fancyhf{}
\lhead{Week 7 / Take-away games / \level}
\lfoot{Bellingham Math Circle / Week 7 / \packetid}
\rfoot{\thepage}
\renewcommand{\headrulewidth}{0pt}
\renewcommand{\footrulewidth}{0pt}

% ---------- problems ----------
\newcounter{prob}
\newlength{\probgap}\setlength{\probgap}{0.35in}
% A problem never splits across pages.
\newenvironment{problem}%
  {\par\noindent\begin{minipage}[t]{\linewidth}\raggedright
   \stepcounter{prob}\textbf{Problem \theprob:}\ }%
  {\end{minipage}\par\vspace{\probgap}}

% ---------- number tracks ----------
% A snake: 0 to 10 run left to right in the first row, 11 sits under 10
% and 11-20 run right to left, 21 sits under 20 and 21-30 run left to
% right. Consecutive squares always share an edge. Thick walls separate
% rows except where the track turns.
\newlength{\sq}\setlength{\sq}{0.676in}
\newcommand{\tracknum}{\sffamily\large}
\newcommand{\track}[1]{%
  \begin{tikzpicture}[x=\sq,y=\sq,line join=miter]
    \foreach \n in {0,...,#1}{%
      \pgfmathtruncatemacro{\r}{\n==0 ? 0 : div(\n-1,10)}%
      \pgfmathtruncatemacro{\k}{\n==0 ? -1 : mod(\n-1,10)}%
      \pgfmathtruncatemacro{\odd}{mod(\r,2)}%
      \pgfmathtruncatemacro{\c}{\odd==1 ? 10-\k : \k+1}%
      \draw[line width=0.8pt] (\c,-\r) rectangle ++(1,1);
      \node[font=\tracknum,inner sep=0pt,text height=1.6ex,text depth=0pt] at (\c+0.5,-\r+0.5) {\n};
    }
    \ifnum#1>20
      \draw[line width=2.2pt] (1,0) -- (10,0);
      \draw[line width=2.2pt] (2,-1) -- (11,-1);
      \draw[line width=2.2pt] (0,1) -- (11,1) -- (11,-2) -- (1,-2) -- (1,0) -- (0,0) -- cycle;
    \else\ifnum#1>10
      \draw[line width=2.2pt] (1,0) -- (10,0);
      \draw[line width=2.2pt] (0,1) -- (11,1) -- (11,-1) -- (1,-1) -- (1,0) -- (0,0) -- cycle;
    \else
      \draw[line width=2.2pt] (0,0) rectangle (11,1);
    \fi\fi
  \end{tikzpicture}}

% ---------- piles of counters ----------
% Use inside a tikzpicture with x=1in,y=1in; arguments are the centre.
\newcommand{\counter}[2]{%
  \draw[line width=1pt,fill=black!22] ({#1},{#2}) circle (0.2in);}

\renewcommand{\tracknum}{\sffamily\Large}
\linespread{1.08}

\newcommand{\gap}{\par\vspace{0.3in}}
\newcommand{\rulelabel}[1]{\par\vspace{0.14in}#1\par\vspace{0.08in}}

% One row per pile: "first" and "second" to circle, and a blank for the
% number of counters taken on the first turn.
\newcommand{\pilerecord}[1]{%
  \begin{tikzpicture}[x=1in,y=1in]
    \foreach \n [count=\i] in {#1}{%
      \pgfmathsetmacro{\y}{-(\i-1)*0.58}
      \node[anchor=base west,font=\Large] at (0.1,\y) {Pile of \n};
      \node[anchor=base,font=\Large] at (2.9,\y) {first};
      \node[anchor=base,font=\Large] at (4.2,\y) {second};
      \node[anchor=base west,font=\Large] at (5.2,\y) {take};
      \draw[line width=0.6pt] (5.85,\y-0.04) -- (6.9,\y-0.04);
    }
  \end{tikzpicture}}

% Chart for two piles: rows are the first pile, columns the second pile.
\newcommand{\pilechart}[1]{%
  \begin{tikzpicture}[x=0.39in,y=0.39in]
    \foreach \a in {0,...,#1}{%
      \foreach \b in {0,...,#1}{\draw[line width=0.9pt] (\b,-\a) rectangle ++(1,-1);}
      \node[font=\sffamily\large,inner sep=0pt,text height=1.6ex,text depth=0pt] at (-0.45,-\a-0.5) {\a};
      \node[font=\sffamily\large,inner sep=0pt,text height=1.6ex,text depth=0pt] at (\a+0.5,0.4) {\a};
    }
    \pgfmathsetmacro{\mid}{(#1+1)/2}
    \node[font=\large] at (\mid,1.1) {second pile};
    \node[font=\large,rotate=90] at (-1.2,-\mid) {first pile};
  \end{tikzpicture}}

\begin{document}

Two players take turns. Whoever takes the last counter wins, and a player
who cannot move while counters are left loses. On a track, the token
starts on the number of counters and moves toward 0 by the number of
counters taken, and whoever puts it on 0 wins. When you play, two of you
are the players and the third is the referee, and you switch roles after
every game.

\vspace{0.25in}

\begin{problem}
Take 1 or 2 counters on each turn. Play each pile with real counters until
you are sure whether you would rather go first or second, and circle your
choice. If you would rather go first, write how many counters you take on
your first turn.

\vspace{0.3in}
\pilerecord{7,9,11,12,16}
\end{problem}

\begin{problem}
For each rule, colour every square where you would like to leave the token
for your opponent.

\rulelabel{Take 1 or 2 counters.}
\track{20}
\rulelabel{Take 1, 2 or 3 counters.}
\track{20}
\end{problem}

\newpage

\begin{problem}
Now take 1, 3 or 4 counters on each turn. Play each pile until you are
sure whether you would rather go first or second, and circle your choice.
If you would rather go first, write how many counters you take on your
first turn.

\vspace{0.3in}
\pilerecord{5,8,9,11,14}
\end{problem}

\begin{problem}
With moves of 1, 3 or 4, colour every square where you would like to leave
the token for your opponent.

\vspace{0.2in}
\track{30}
\end{problem}

\newpage

\begin{problem}
Kai plays with moves of 1, 3 or 4. He goes first and takes as many counters
as he is allowed on his first turn. Find every pile from 1 to 30 where the
player who goes first can win, but Kai's first move leaves a pile from
which his opponent can win. Colour those piles.

\vspace{0.2in}
\track{30}
\end{problem}

\begin{problem}
A pile has 100 counters. Would you rather go first or second with moves of
1 or 2? With moves of 1, 2 or 3? With moves of 1, 3 or 4? Whenever you
would go first, write how many counters you take. Explain how you would
win whatever your opponent does.
\vspace{3.2in}
\end{problem}

\newpage

\begin{problem}
For each rule, colour every square where you would like to leave the token
for your opponent.

\rulelabel{Take 1 or 4 counters.}
\track{30}
\rulelabel{Take 2 or 3 counters.}
\track{30}
\rulelabel{Take 1, 3 or 5 counters.}
\track{30}
\end{problem}

\newpage

\begin{problem}
With moves of 1, 3 or 4, does the pattern of coloured squares keep going
the same way forever? Explain.
\vspace{3in}
\end{problem}

\begin{problem}
Choose which moves are allowed so that the squares you would like to leave
are exactly 0, 5, 10, 15, and so on. Then find as many different choices as
you can that make them exactly 0, 3, 6, 9, and so on. Explain how you can
tell which choices work.
\vspace{3.5in}
\end{problem}

\newpage

\begin{problem}
Now whoever takes the last counter loses. Take 1 or 2 counters on each
turn. Play each pile until you are sure whether you would rather go first
or second, and circle your choice. If you would rather go first, write how
many counters you take on your first turn.

\vspace{0.3in}
\pilerecord{1,3,5,7,8}
\end{problem}

\begin{problem}
Whoever takes the last counter still loses, so whoever puts the token on 0
loses. For each rule, colour every square where you would like to leave the
token for your opponent.

\rulelabel{Take 1 or 2 counters.}
\track{20}
\rulelabel{Take 1, 2 or 3 counters.}
\track{20}
\end{problem}

\newpage

\begin{problem}
Now there are two piles. On your turn, take 1 or 2 counters from one pile,
whichever you choose. Whoever takes the last counter on the table wins.
Play piles of 7 and 7. Would you rather go first or second? Explain why
your way of playing wins whatever your opponent does.

\vspace{0.35in}
\begin{tikzpicture}[x=1in,y=1in]
  \foreach \k in {1,...,7}{\counter{(\k-1)*0.47+0.25}{0}}
  \foreach \k in {1,...,7}{\counter{(\k-1)*0.47+4.0}{0}}
\end{tikzpicture}
\vspace{1.75in}
\end{problem}

\begin{problem}
In the two-pile game of Problem 12, colour every square of the chart where
you would like to leave the two piles for your opponent. Then decide
whether you would rather go first or second with piles of 20 and 11, and
explain.

\vspace{0.15in}
\hspace*{0.5in}\pilechart{9}
\end{problem}

\newpage

\begin{problem}
Kim says that for any list of allowed moves, such as 2, 5 or 6, the coloured
squares on a long enough track end up repeating. Is Kim right? Explain.

\rulelabel{Take 2, 5 or 6 counters.}
\track{30}
\end{problem}

\end{document}
